\section{Problem definition}
\label{sec:problem-definition}

Let \(G=(V,A,w)\) be a directed graph with vertex set \(V\), arc set \(A\subseteq V\times V\), and nonnegative arc weights \(w:A\to \mathbb{R}_{\ge 0}\). A feedback arc set is a set \(F\subseteq A\) such that removing \(F\) makes the graph acyclic. The minimum weighted feedback arc set problem asks for a minimum-weight feedback arc set, i.e., a set \(F\) minimizing \(\sum_{a\in F}w(a)\).

An ordering \(\pi\) assigns each vertex a distinct position. An arc \((u,v)\) is backward under \(\pi\) if \(u\) appears after \(v\). The ordering-induced backward weight is

\begin{equation}
    \operatorname{bw}(\pi)
    =
    \sum_{(u,v)\in A:\ \pi(u)>\pi(v)} w(u,v).
    \label{eq:backward-weight}
\end{equation}

For any ordering \(\pi\), the set of backward arcs is a feedback arc set: all remaining forward arcs point from earlier to later positions and therefore form an acyclic graph. Conversely, every acyclic subgraph admits a topological order. The weighted feedback arc set problem can therefore be treated as the problem of finding an ordering with minimum backward weight \cite{F90,K72}.

Many constructive heuristics in this study first build an acyclic \emph{active} subgraph \(H=(V,A\setminus F)\) by deleting a set \(F\subseteq A\), and then extract a linear extension \(\pi\) of \(H\). For such a pair \((F,\pi)\), define the backward arc set induced by \(\pi\) on the original graph as
\begin{equation}
    B_{\pi}
    =
    \{(u,v)\in A:\ \pi(v)<\pi(u)\}.
    \label{eq:backward-set}
\end{equation}
Every active arc is forward under \(\pi\), so \(B_{\pi}\subseteq F\). For nonnegative weights,
\begin{equation}
    \operatorname{bw}(\pi)=w(B_{\pi})\le w(F).
    \label{eq:fas-order-inequality}
\end{equation}
Equality holds for a given \(\pi\) if and only if every removed arc that carries positive weight is backward under \(\pi\); equivalently, \(w(F\setminus B_{\pi})=0\). If all removed arcs have strictly positive weight, then \(\operatorname{bw}(\pi)=w(F)\) if and only if \(B_{\pi}=F\). Topological extraction is therefore \emph{not} objective-neutral in general: different linear extensions of the same active DAG can yield different values of \(\operatorname{bw}(\pi)\). The implementation reports the ordering objective \(\operatorname{bw}(\pi)\) recomputed from the returned ranking, not the removed-set weight \(w(F)\) alone.

Dense linear ordering problem benchmarks are often stated in terms of a complete weight matrix \(C\), where an ordering \(\pi=(\pi_1,\ldots,\pi_n)\) is evaluated by the forward weight. In that notation,

\begin{equation}
    \operatorname{fw}_{C}(\pi)
    =
    \sum_{1\le i<j\le n} C_{\pi_i,\pi_j},
    \qquad
    \operatorname{bw}_{C}(\pi)
    =
    \sum_{1\le i<j\le n} C_{\pi_j,\pi_i}
    =
    W_C-\operatorname{fw}_{C}(\pi),
    \label{eq:lolib-forward-backward}
\end{equation}

where \(W_C=\sum_{i\ne j}C_{ij}\) is the total off-diagonal weight. Thus maximizing forward weight on a complete dense linear-ordering instance is equivalent to minimizing backward weight after converting the matrix to a complete weighted digraph. LOLIB results are reported as a transfer test. Those instances are dense complete ordering problems rather than sparse directed graph benchmarks.
